% SHARD-ID: CA-48-QUBIT-MOMENT-MATRIX-CONSTRUCTION
% SHARD-TITLE: Moment Matrices and Positivity Constraints
% SHARD-SUMMARY: Builds the finite moment matrix M_{u,v}=omega(P_u^*P_v) from exact Pauli products.
% SHARD-SUMMARY: Uses the realification [A -B; B A] to impose complex Hermitian positivity in JuMP.
% SHARD-SUMMARY: Checks feasible and infeasible tiny Mosek instances as smoke tests.
% SHARD-KEYWORDS: moment matrix, positivity, realification, JuMP, Mosek, qubit SDP

\section{Moment Matrices and Positivity Constraints}
\label{sec:qubit-moment-matrix-construction}

\subsection{Status, Sources, and Dependencies}

\textbf{Status: checked.}  Moment matrices are constructed by
\texttt{build\_qubit\_sdp\_model}.

Dependencies: CA-43--CA-47.

% Source: literature/md/2010.11121/2010.11121.md:178--184 -- projectively
%   consistent states and GNS compatibility.
% Checked: src/QubitMomentSDP.jl; test/runtests.jl testset
%   "qubit Pauli moment SDP hierarchy".

\subsection{Affine Moment Matrix}

For a finite actual probe set \(\mathcal W\), define
\[
  M_{u,v}
  =
  \omega(P_u^*P_v)
  =
  \phi(u,v)y_{\operatorname{can}(u\cdot v)}.
\]
The identity moment is the constant \(1\).  All other canonical words appearing
in products or relation constraints receive real JuMP variables.

Write \(M=A+iB\), with \(A,B\) real affine matrices.  Positivity of the complex
Hermitian moment matrix is imposed by the real symmetric cone
\[
  \begin{pmatrix}
  A&-B\\
  B&A
  \end{pmatrix}\succeq0.
  \tag{48.1}
\]
For the one-site probe set \(\{I,X,Y,Z\}\), the implementation creates three
real moment variables and an \(8\times8\) realified PSD block.  This exact
dimension is checked in the Julia suite.

\subsection{Tiny Feasibility Tests}

The zero-residual one-site moment problem solves with Mosek as
\[
  \texttt{:not\_excluded\_at\_level}.
\]
The artificial relation \(\omega(I)=0\) solves as \(\texttt{:excluded}\).
These are smoke tests for the solver path, not physical Hamiltonian claims.
